% GENERATED FILE. Edit canonical lessons, then run npm run generate:books.
% canonical-source-hash: fnv1a64:b28105336169077e
% canonical-lessons: SA-C149-karmakarah, SA-C149-sramikah, SA-C149-adhikari, SA-C149-lipikah, SA-C149-nayakah, SA-R149-first-pass-words-from-chapters-141-144, SA-R149-second-pass-words-from-chapters-145-149

\chapter{Worker, Labourer, Official, Clerk, Leader}
\label{ch:sa-worker-labourer-official-clerk-leader}
\hlchaptermodality{149}

\begin{chapteropening}
\textbf{By the end of this chapter:} \emph{I can say a worker, a labourer, an official, a clerk and a leader.}

Complete the last lesson of chapter 149: I can say a worker, a labourer, an official, a clerk and a leader.
\end{chapteropening}

\section[\texorpdfstring{\sk{कर्मकरः} (karmakaraḥ)}{कर्मकरः (karmakaraḥ)}]{\sk{कर्मकरः} (karmakaraḥ) — a worker, a labourer}

\label{lesson:SA-C149-karmakarah}



\begin{quote}
Before the new one: say the Sanskrit for a fisherman, then the Sanskrit for a cowherd.
\end{quote}

\subsection*{You'll want to know: \sk{कर्मकरः}}
\textbf{\sk{कर्मकरः}} — \emph{karmakaraḥ} — \textquotedblleft{}a worker, a labourer\textquotedblright{}.

\begin{grammarlens}[title={What you've built}]
One more word for this chapter.
\end{grammarlens}

\subsection*{Your turn}
\begin{itemize}
\raggedright
  \item \emph{Say it:} \emph{karmakaraḥ}
  \item \emph{Say it:} \emph{karmakaraḥ}, once more
  \item \emph{Say it:} say \emph{karmakaraḥ}
  \item \emph{Recall:} say the Sanskrit for a fisherman, then the Sanskrit for a cowherd, then say \emph{karmakaraḥ} again
\end{itemize}

\begin{tcolorbox}[breakable,colback=teal!4,colframe=teal!35!black,title={Before you move on}]
What does \sk{कर्मकरः} mean? (\textquotedblleft{}A worker, a labourer\textquotedblright{}.) Say it once more.
\end{tcolorbox}
\section[\texorpdfstring{\sk{श्रमिकः} (śramikaḥ)}{श्रमिकः (śramikaḥ)}]{\sk{श्रमिकः} (śramikaḥ) — a labourer}

\label{lesson:SA-C149-sramikah}



\begin{quote}
Before the new one: say the Sanskrit for a cowherd, then the Sanskrit for a worker.
\end{quote}

\subsection*{You'll want to know: \sk{श्रमिकः}}
\textbf{\sk{श्रमिकः}} — \emph{śramikaḥ} — \textquotedblleft{}a labourer\textquotedblright{}.

\begin{grammarlens}[title={What you've built}]
One more word for this chapter.
\end{grammarlens}

\subsection*{Your turn}
\begin{itemize}
\raggedright
  \item \emph{Say it:} \emph{śramikaḥ}
  \item \emph{Say it:} \emph{śramikaḥ}, once more
  \item \emph{Say it:} say \emph{śramikaḥ}
  \item \emph{Recall:} say the Sanskrit for a cowherd, then the Sanskrit for a worker, then say \emph{śramikaḥ} again
\end{itemize}

\begin{tcolorbox}[breakable,colback=teal!4,colframe=teal!35!black,title={Before you move on}]
What does \sk{श्रमिकः} mean? (\textquotedblleft{}A labourer\textquotedblright{}.) Say it once more.
\end{tcolorbox}
\section[\texorpdfstring{\sk{अधिकारी} (adhikārī)}{अधिकारी (adhikārī)}]{\sk{अधिकारी} (adhikārī) — an official}

\label{lesson:SA-C149-adhikari}



\begin{quote}
Before the new one: say the Sanskrit for a worker, then the Sanskrit for a labourer.
\end{quote}

\subsection*{You'll want to know: \sk{अधिकारी}}
\textbf{\sk{अधिकारी}} — \emph{adhikārī} — \textquotedblleft{}an official\textquotedblright{}.

\begin{grammarlens}[title={What you've built}]
One more word for this chapter.
\end{grammarlens}

\subsection*{Your turn}
\begin{itemize}
\raggedright
  \item \emph{Say it:} \emph{adhikārī}
  \item \emph{Say it:} \emph{adhikārī}, once more
  \item \emph{Say it:} say \emph{adhikārī}
  \item \emph{Recall:} say the Sanskrit for a worker, then the Sanskrit for a labourer, then say \emph{adhikārī} again
\end{itemize}

\begin{tcolorbox}[breakable,colback=teal!4,colframe=teal!35!black,title={Before you move on}]
What does \sk{अधिकारी} mean? (\textquotedblleft{}An official\textquotedblright{}.) Say it once more.
\end{tcolorbox}
\section[\texorpdfstring{\sk{लिपिकः} (lipikaḥ)}{लिपिकः (lipikaḥ)}]{\sk{लिपिकः} (lipikaḥ) — a clerk}

\label{lesson:SA-C149-lipikah}



\begin{quote}
Before the new one: say the Sanskrit for a labourer, then the Sanskrit for an official.
\end{quote}

\subsection*{You'll want to know: \sk{लिपिकः}}
\textbf{\sk{लिपिकः}} — \emph{lipikaḥ} — \textquotedblleft{}a clerk\textquotedblright{}.

\begin{grammarlens}[title={What you've built}]
One more word for this chapter.
\end{grammarlens}

\subsection*{Your turn}
\begin{itemize}
\raggedright
  \item \emph{Say it:} \emph{lipikaḥ}
  \item \emph{Say it:} \emph{lipikaḥ}, once more
  \item \emph{Say it:} say \emph{lipikaḥ}
  \item \emph{Recall:} say the Sanskrit for a labourer, then the Sanskrit for an official, then say \emph{lipikaḥ} again
\end{itemize}

\begin{tcolorbox}[breakable,colback=teal!4,colframe=teal!35!black,title={Before you move on}]
What does \sk{लिपिकः} mean? (\textquotedblleft{}A clerk\textquotedblright{}.) Say it once more.
\end{tcolorbox}
\section[\texorpdfstring{\sk{नायकः} (nāyakaḥ)}{नायकः (nāyakaḥ)}]{\sk{नायकः} (nāyakaḥ) — a leader}

\label{lesson:SA-C149-nayakah}



\begin{quote}
Before the new one: say the Sanskrit for an official, then the Sanskrit for a clerk.
\end{quote}

\subsection*{You'll want to know: \sk{नायकः}}
\textbf{\sk{नायकः}} — \emph{nāyakaḥ} — \textquotedblleft{}a leader\textquotedblright{}.

\begin{grammarlens}[title={What you've built}]
One more word for this chapter.
\end{grammarlens}

\subsection*{Your turn}
\begin{itemize}
\raggedright
  \item \emph{Say it:} \emph{nāyakaḥ}
  \item \emph{Say it:} \emph{nāyakaḥ}, once more
  \item \emph{Say it:} say \emph{nāyakaḥ}
  \item \emph{Recall:} say the Sanskrit for an official, then the Sanskrit for a clerk, then say \emph{nāyakaḥ} again
\end{itemize}

\begin{tcolorbox}[breakable,colback=teal!4,colframe=teal!35!black,title={Before you move on}]
What does \sk{नायकः} mean? (\textquotedblleft{}A leader\textquotedblright{}.) Say it once more.
\end{tcolorbox}
\section[(dialogue)]{Words from chapters 141-144 — a first pass}

\label{lesson:SA-R149-first-pass-words-from-chapters-141-144}



\begin{quote}
Say the words for a clerk and a leader.
\end{quote}

\subsection*{Your turn}
Say these aloud:

\begin{itemize}
\raggedright
  \item drinks, dwells, moves about, touches, leaves
  \item is pleased, gets tired, composes, shows, worships
  \item comes back, lies down, sits, is afraid, meditates
  \item shares, weighs, flies up, puts in, joins
\end{itemize}

\begin{tcolorbox}[breakable,colback=teal!4,colframe=teal!35!black,title={Before you move on}]
Drinks? (\textbf{\sk{पिबति}}, \emph{pibati}.) Dwells? (\textbf{\sk{वसति}}, \emph{vasati}.) Moves about? (\textbf{\sk{चरति}}, \emph{carati}.) Touches? (\textbf{\sk{स्पृशति}}, \emph{spṛśati}.) Leaves? (\textbf{\sk{त्यजति}}, \emph{tyajati}.) Is pleased? (\textbf{\sk{तुष्यति}}, \emph{tuṣyati}.) Gets tired? (\textbf{\sk{श्राम्यति}}, \emph{śrāmyati}.) Composes? (\textbf{\sk{रचयति}}, \emph{racayati}.) Shows? (\textbf{\sk{दर्शयति}}, \emph{darśayati}.) Worships? (\textbf{\sk{पूजयति}}, \emph{pūjayati}.) Comes back? (\textbf{\sk{प्रत्यागच्छति}}, \emph{pratyāgacchati}.) Lies down? (\textbf{\sk{शेते}}, \emph{śete}.) Sits? (\textbf{\sk{आस्ते}}, \emph{āste}.) Is afraid? (\textbf{\sk{बिभेति}}, \emph{bibheti}.) Meditates? (\textbf{\sk{ध्यायति}}, \emph{dhyāyati}.) Shares? (\textbf{\sk{भजति}}, \emph{bhajati}.) Weighs? (\textbf{\sk{तोलयति}}, \emph{tolayati}.) Flies up? (\textbf{\sk{उत्पतति}}, \emph{utpatati}.) Puts in? (\textbf{\sk{निक्षिपति}}, \emph{nikṣipati}.) Joins? (\textbf{\sk{योजयति}}, \emph{yojayati}.)
\end{tcolorbox}
\section[(dialogue)]{Words from chapters 145-149 — a second pass}

\label{lesson:SA-R149-second-pass-words-from-chapters-145-149}



\begin{quote}
Say the words for is glad and cries out.
\end{quote}

\subsection*{Your turn}
Say these aloud:

\begin{itemize}
\raggedright
  \item is glad, cries out, roars, coos, chatters
  \item a grandson, a granddaughter, a paternal grandmother, a maternal grandfather, a father's brother
  \item a master, a servant, a queen, a minister, a general
  \item a barber, a cook, a gardener, a fisherman, a cowherd
  \item a worker, a labourer, an official, a clerk, a leader
\end{itemize}

\begin{tcolorbox}[breakable,colback=teal!4,colframe=teal!35!black,title={Before you move on}]
Is glad? (\textbf{\sk{नन्दति}}, \emph{nandati}.) Cries out? (\textbf{\sk{क्रन्दति}}, \emph{krandati}.) Roars? (\textbf{\sk{गर्जति}}, \emph{garjati}.) Coos? (\textbf{\sk{कूजति}}, \emph{kūjati}.) Chatters? (\textbf{\sk{जल्पति}}, \emph{jalpati}.) A grandson? (\textbf{\sk{पौत्रः}}, \emph{pautraḥ}.) A granddaughter? (\textbf{\sk{पौत्री}}, \emph{pautrī}.) A paternal grandmother? (\textbf{\sk{पितामही}}, \emph{pitāmahī}.) A maternal grandfather? (\textbf{\sk{मातामहः}}, \emph{mātāmahaḥ}.) A father's brother? (\textbf{\sk{पितृव्यः}}, \emph{pitṛvyaḥ}.) A master? (\textbf{\sk{स्वामी}}, \emph{svāmī}.) A servant? (\textbf{\sk{सेवकः}}, \emph{sevakaḥ}.) A queen? (\textbf{\sk{राज्ञी}}, \emph{rājñī}.) A minister? (\textbf{\sk{मन्त्री}}, \emph{mantrī}.) A general? (\textbf{\sk{सेनापतिः}}, \emph{senāpatiḥ}.) A barber? (\textbf{\sk{नापितः}}, \emph{nāpitaḥ}.) A cook? (\textbf{\sk{पाचकः}}, \emph{pācakaḥ}.) A gardener? (\textbf{\sk{मालाकारः}}, \emph{mālākāraḥ}.) A fisherman? (\textbf{\sk{धीवरः}}, \emph{dhīvaraḥ}.) A cowherd? (\textbf{\sk{गोपालः}}, \emph{gopālaḥ}.) A worker? (\textbf{\sk{कर्मकरः}}, \emph{karmakaraḥ}.) A labourer? (\textbf{\sk{श्रमिकः}}, \emph{śramikaḥ}.) An official? (\textbf{\sk{अधिकारी}}, \emph{adhikārī}.) A clerk? (\textbf{\sk{लिपिकः}}, \emph{lipikaḥ}.) A leader? (\textbf{\sk{नायकः}}, \emph{nāyakaḥ}.)
\end{tcolorbox}
